\section{Sharp limitations of the curvature-coordinate comparison}
\label{sec:coordinate}
The coordinate $h=\int\sqrt g$ is classical in separable Bregman
geometry~\cite{gomes,gzyl}. The next uniform estimate is the Phase~8
comparison~\cite{huang8}, reproduced to make the obstruction self-contained.
Let $g=\varphi''$ be a nonzero nonnegative polynomial on $[0,1]$
of degree at most $r\ge1$. Put $K_r=2\log(64r^2)+1/4$ and
$h(x)=\int_0^x\sqrt{g(t)}\,dt$.

\begin{lemma}[Polynomial curvature-coordinate comparison]
\label{lem:poly-h-new}
For every $x,y\in[0,1]$ and both orders of the arguments,
\begin{equation}
 \frac{1}{K_r}|h(x)-h(y)|^2
 \le B_\varphi(x,y)
 \le 8r^2|h(x)-h(y)|^2.
 \label{eq:poly-h-new}
\end{equation}
The constants are uniform over the polynomial coefficients and over the
location and length of the interval between $x$ and $y$.
\end{lemma}

\begin{proof}
The assertion is immediate for $x=y$. Fix $u<v$, and put
$L=v-u$, $\mu=\int_u^v g(t)\,\mathrm dt$, and
$M=\max_{[u,v]}g$. Both $\mu$ and $M$ are positive.
The Markov mass estimates used here are
\begin{equation}
 M\le\frac{8r^2\mu}{L},\qquad
 \int_u^v|t-a|g(t)\,\mathrm dt\ge\frac{L\mu}{64r^2}
 \quad (a\in[u,v]).
 \label{eq:poly-markov-new}
\end{equation}
For completeness, Markov's classical polynomial inequality
\cite{shadrin}, rescaled to $[u,v]$, gives
$\|g'\|_\infty\le2r^2M/L$. From a maximizer there is an available
one-sided segment of length $L/(4r^2)$, on which $g\ge M/2$.
Integration gives the first estimate. For any $a\in[u,v]$, the mass in
$\{t\in[u,v]:|t-a|<L/(32r^2)\}$ is at most
$LM/(16r^2)\le\mu/2$. The complementary mass gives the second estimate.

The two divergence orientations are
\[
 B_\varphi(v,u)=\int_u^v(v-t)g(t)\,\mathrm dt,
 \qquad
 B_\varphi(u,v)=\int_u^v(t-u)g(t)\,\mathrm dt.
\]
Fix either of these values and denote it by $B$. Measure the coordinate
$s\in[0,L]$ from the endpoint where its linear kernel vanishes, and denote
the resulting curvature by $\widetilde g(s)$. Explicitly,
$\widetilde g(s)=g(v-s)$ for $B_\varphi(v,u)$, and
$\widetilde g(s)=g(u+s)$ for $B_\varphi(u,v)$. Then
\[
 B=\int_0^L s\widetilde g(s)\,\mathrm ds,
 \qquad A:=h(v)-h(u)=\int_0^L\sqrt{\widetilde g(s)}\,\mathrm ds.
\]
Taking $a$ to be the zero-kernel endpoint in
\eqref{eq:poly-markov-new} gives
\begin{equation}
 \frac{L\mu}{64r^2}\le B\le L\mu,
 \qquad
 M\le\frac{512r^4B}{L^2}.
 \label{eq:poly-BM-new}
\end{equation}

Set $\delta=L/(64r^2)\in(0,L)$ and split $A=A_0+A_1$ at
$s=\delta$. The short endpoint contribution satisfies
\[
 A_0^2=\left(\int_0^\delta\sqrt{\widetilde g(s)}\,\mathrm ds\right)^2
 \le\delta^2M\le\frac{B}{8}.
\]
Weighted Cauchy--Schwarz gives
\[
 A_1^2
 =\left(\int_\delta^L
       \sqrt{s\widetilde g(s)}\,s^{-1/2}\,\mathrm ds\right)^2
 \le\left(\int_\delta^L s\widetilde g(s)\,\mathrm ds\right)
       \left(\int_\delta^L\frac{\mathrm ds}{s}\right)
 \le B\log(64r^2).
\]
Therefore $A^2\le2A_0^2+2A_1^2\le K_rB$, which proves the
left inequality in \eqref{eq:poly-h-new}.

For the other direction, $\sqrt{g(t)}\ge g(t)/\sqrt M$ implies
\[
 A^2\ge\frac{\mu^2}{M}\ge\frac{L\mu}{8r^2}.
\]
Together with $B\le L\mu$, this proves $B\le8r^2A^2$.
Every estimate applies to either orientation, completing the proof.
\end{proof}


For later use, define the transformed partition variance by
\[
 V_h(P)=\sum_{C\in P}\sum_{i\in C}p_i
 \bigl(h(x_i)-\overline h_C\bigr)^2,\qquad
 \overline h_C=p_C^{-1}\sum_{i\in C}p_i h(x_i).
\]
Minimizing the two sides of~\eqref{eq:poly-h-new} over a center in
each cell gives $K_r^{-1}V_h(P)\le D_\varphi(P)\le8r^2V_h(P)$.
This uses that $h([0,1])$ is an interval, and does not assume
$h(m_C)=\overline h_C$.

\subsection{Separate and simultaneous obstructions}

Let $g\geq0$ be a nonzero polynomial on $[0,1]$, let
$\varphi''=g$, and put
\begin{align*}
 h(x)&=\int_0^x\sqrt{g(t)}\,dt,&
 A&=h(1)-h(0),\\
 B_-&=B_\varphi(0,1)=\int_0^1 t g(t)\,dt,&
 B_+&=B_\varphi(1,0)=\int_0^1(1-t)g(t)\,dt.
\end{align*}
The orientation in these definitions is intentional. Define the two
degree-uniform comparison constants
\[
 L_r=\sup_{\substack{g\geq0,\ g\not\equiv0\\\deg g\leq r}}
       \ \sup_{u\ne v}\frac{|h(u)-h(v)|^2}{B_\varphi(u,v)},
 \qquad
 U_r=\sup_{\substack{g\geq0,\ g\not\equiv0\\\deg g\leq r}}
       \ \sup_{u\ne v}\frac{B_\varphi(u,v)}{|h(u)-h(v)|^2},
\]
where $u,v\in[0,1]$ and each occurrence of $h$ and $\varphi$ belongs to
the same $g$ inside its supremum.

\begin{proposition}[The inverse logarithm cannot be removed]
For every integer $m\geq2$, there is strictly positive polynomial curvature
of degree $2m-2$ for which
\[
 \frac{A^2}{B_-}=\frac{H_m^2}{2H_m-H_{2m}}\sim\log m,
 \qquad H_m=\sum_{j=1}^m\frac1j.
\]
Consequently $L_r=\Theta(\log(r+2))$ by Lemma~\ref{lem:poly-h-new}. In particular no absolute $C$ can give
$|h(u)-h(v)|^2\leq C B_\varphi(u,v)$ for all polynomial degrees.
\end{proposition}
\begin{proof}
Set
\[
 q_m(t)=\sum_{j=0}^{m-1}(1-t)^j
       =\frac{1-(1-t)^m}{t},\qquad g_m(t)=q_m(t)^2,
\]
with the polynomial continuation $q_m(0)=m$. On $[0,1]$, $q_m\geq1$,
so $\sqrt{g_m}=q_m$ and $A=H_m$. For every integer $k\geq1$,
\[
 \int_0^1\frac{1-(1-t)^k}{t}\,dt
 =\sum_{j=0}^{k-1}\int_0^1(1-t)^j\,dt=H_k.
\]
The identity
$[1-(1-t)^m]^2=2[1-(1-t)^m]-[1-(1-t)^{2m}]$
therefore yields $B_-=2H_m-H_{2m}$. Finally
$H_m=\log m+O(1)$ and $H_{2m}-H_m=\log2+o(1)$, proving the
asymptotic. Taking $m=\lfloor r/2\rfloor+1$ supplies a witness for every
degree-at-most-$r$ envelope, rather than only an even subsequence.
\end{proof}

\begin{proposition}[The forward quadratic loss cannot gain a logarithm]
For every integer $m\geq2$, there is nonnegative nonzero polynomial
curvature of degree $4m-4$ such that
\[
 \frac{B_+}{A^2}\geq\frac{m^2}{32768}.
\]
Consequently $U_r=\Theta((r+1)^2)$ by Lemma~\ref{lem:poly-h-new}. In particular an absolute bound
$B_\varphi(u,v)\leq C r^2\,|h(u)-h(v)|^2/\log(r+2)$ is false.
The same asymptotic obstruction holds if curvature is required to be
strictly positive.
\end{proposition}
\begin{proof}
Let $T_m$ be the Chebyshev polynomial characterized by
$T_m(\cos\theta)=\cos(m\theta)$. The quotient
\[
 F_m(t)=\frac{1-T_m(1-2t)}{2m^2t},\qquad F_m(0)=1,
\]
is a polynomial of degree $m-1$. Writing $\theta=\arcsin\sqrt t$
shows that
\[
 F_m(t)=\left(\frac{\sin(m\theta)}{m\sin\theta}\right)^2,
 \qquad 0\leq F_m(t)\leq
     \min\left\{1,\frac1{m^2t}\right\}.
\]
Here $|\sin(m\theta)|\leq m\sin\theta$ follows from the geometric-sum
formula for $\sin(m\theta)/\sin\theta$, and the other bound uses
$|\sin(m\theta)|\leq1$.

Set $g_m=F_m^4$, of degree $4m-4$. Since $F_m\geq0$,
\[
 A=\int_0^1F_m(t)^2\,dt
 \leq\frac1{m^2}+\int_{1/m^2}^1\frac{dt}{m^4t^2}
 =\frac2{m^2}-\frac1{m^4}\leq\frac2{m^2}.
\]
If $0\leq t\leq1/(16m^2)$, then
$\sqrt t\leq\theta\leq2\sqrt t$ and $m\theta\leq1/2$.
The elementary estimate $\sin y\geq y/2$ for $0\leq y\leq1/2$
therefore gives $F_m(t)\geq1/4$. On this same interval $1-t\geq1/2$,
whence
\[
 B_+\geq\int_0^{1/(16m^2)}(1-t)F_m(t)^4\,dt
 \geq\frac1{8192m^2}.
\]
Division by $A^2\leq4/m^4$ proves the claimed bound. Taking
$m=\lfloor r/4\rfloor+1$ covers every sufficiently large
degree-at-most-$r$ envelope.

For strict positivity one may use $\widetilde g_m=F_m^4+m^{-8}$,
without changing the degree. Since $\sqrt{a+b}\leq\sqrt a+\sqrt b$,
its coordinate length is at most $2/m^2+m^{-4}\leq3/m^2$, whereas its
$B_+$ is at least $1/(8192m^2)$. Thus its ratio is at least
$m^2/73728$. Both examples can be rescaled to have curvature at most one;
positive curvature scaling changes $B$ and $A^2$ by the same factor.
\end{proof}

\begin{proposition}[Both losses occur for one generator and one source]
\label{prop:coordinate-same-source}
There are strictly positive polynomial curvatures of degree $r=4m-4$
for which the product of the best two comparison constants in the
coordinate $h=\int\sqrt g$ is $\Omega(r^2\log r)$. This remains true
for the best two-sided partition-cost comparison constants chosen for a
single three-point source. Nevertheless that source has deterministic
and stochastic all-capacity hierarchy price exactly one.
\end{proposition}
\begin{proof}
Using the two polynomials above, set
\[
 G_m(t)=m^4F_m(t)^4+q_m(1-t)^2.
\]
It is strictly positive on $[0,1]$ and has degree $4m-4$ for $m\geq2$;
when $m=2$ both terms have degree four or two, and in every case the
higher-degree coefficient is positive. Write $h'=\sqrt{G_m}$.

On the left half interval, the inequality
$\sqrt{a+b}\leq\sqrt a+\sqrt b$ and the preceding estimates imply
\[
 A_L:=h(1/2)-h(0)
 \leq m^2\int_0^1F_m^2+\int_0^{1/2}\frac{dt}{1-t}
 \leq2+\log2<3.
\]
The lower estimate $F_m\geq1/4$ on $[0,1/(16m^2)]$ gives
\[
 B_\varphi(1/2,0)
 =\int_0^{1/2}(1/2-t)G_m(t)\,dt
 \geq\frac{m^2}{16384},\qquad
 \frac{B_\varphi(1/2,0)}{A_L^2}\geq\frac{m^2}{147456}.
\]
On the right half, change variables $s=1-t$. Then
\[
 A_R:=h(1)-h(1/2)
 \geq\int_0^{1/2}q_m(s)\,ds
 =H_m-\sum_{j=1}^m\frac{2^{-j}}j\geq H_m-\log2.
\]
For $t\geq1/2$ we have $m^4F_m(t)^4\leq16/m^4$, so
\[
 \begin{split}
 B_\varphi(1,1/2)
 &=\int_{1/2}^1(1-t)G_m(t)\,dt\\
 &\leq 2H_m-H_{2m}+\frac2{m^4}
 \leq H_m+\frac2{m^4}.
 \end{split}
\]
Consequently
\[
 \frac{A_R^2}{B_\varphi(1,1/2)}
 \geq\frac{(H_m-\log2)^2}{H_m+2/m^4}\sim\log m.
\]
The two ratios for this same generator therefore have product
$\Omega(m^2\log m)$.

To transfer this obstruction to partition comparisons, consider the
three locations $0,1/2,1$ with weights proportional to
$1,\eta,\eta^2$, where $\eta>0$. Let $P_L$ merge the first two
locations and let $P_R$ merge the last two. The other location is a
singleton in each partition. In general a two-point cell with locations
$x,y$ and weights $a,\eta a$ satisfies, as $\eta\downarrow0$,
\[
 \frac{D_\varphi(\{x,y\})}{a\eta}\longrightarrow B_\varphi(y,x),
 \qquad
 \frac{V_h(\{x,y\})}{a\eta}\longrightarrow
       |h(y)-h(x)|^2.
\]
The first limit follows by differentiating
$a\varphi(x)+a\eta\varphi(y)-a(1+\eta)
\varphi((x+\eta y)/(1+\eta))$ at $\eta=0$; the second follows from
the exact weighted two-point variance. Overall normalization cancels
from their quotient. It follows that
\[
 \frac{D_\varphi(P_L)}{V_h(P_L)}\longrightarrow
      \frac{B_\varphi(1/2,0)}{A_L^2},\qquad
 \frac{V_h(P_R)}{D_\varphi(P_R)}\longrightarrow
      \frac{A_R^2}{B_\varphi(1,1/2)}.
\]
For each $m$, choose sufficiently small positive $\eta$ so both ratios
are at least half their limits. Thus any constants $c,C>0$ valid for
\emph{every partition of this source},
$cV_h(P)\leq D_\varphi(P)\leq C V_h(P)$, necessarily satisfy
$C/c=\Omega(m^2\log m)$. This choice permits $\eta$ to depend on $m$;
no assertion of a fixed minimum source mass is made.

Finally, for three states the single nontrivial capacity is two. Choose
its flat optimum and extend by the singleton partition at capacity
three and the one-cell partition at capacity one. This is a hierarchy
with price one. Stochastic price is also one because its flat
denominators coincide with the deterministic ones.
\end{proof}

\paragraph{Exact scope of the obstruction.}
The separate orders of the comparison constants are sharp, and
their product is sharp even when both constants can depend on the
generator and source. Thus a transfer using a global two-sided
comparison of all partition costs in this particular coordinate cannot
improve its $O(r^2\log r)$ distortion factor in general. A method that
uses additional structure or a different coordinate is not excluded.
These statements do not give any growing hierarchy lower bound: the
last construction deliberately has hierarchy price one. Divergence or
partition-comparison distortion and hierarchy incompatibility are
distinct claims.
